\documentclass[10pt, letterpaper]{article}
% ----------------------------------------------------------------------
% LaTeX Preamble
% ----------------------------------------------------------------------

% --- Required Packages ---
\usepackage{reddy_resume}   % Custom resume style package

% ----------------------------------------------------------------------
% Document Body
% ----------------------------------------------------------------------
\begin{document}

% ----------------------------------------------------------------------
% HEADER SECTION
% ----------------------------------------------------------------------
\begin{header}
	% --- Left Column: Name and Contact Info ---
	\begin{minipage}[t]{0.65\textwidth}
		\fontsize{20pt}{20pt}\selectfont \textbf{Aadarsha Gopala Reddy}\\
		\normalsize
		% Address (Optional - Uncomment if needed)
		% \mbox{St. Louis, Missouri 63112 USA}%
		% \kern \kernSpacing%
		% \AND%
		% \kern \kernSpacing%
		% Email
		\mbox{\hrefWithoutArrow{mailto:adurs2002@gmail.com}{adurs2002@gmail.com}}%
		\kern \kernSpacing%
		\AND%
		\kern \kernSpacing%
		% Phone
		\mbox{\hrefWithoutArrow{tel:+17408021776}{+1 (740) 802-1776}}%
		% Website Link (Optional - Uncomment if preferred in left column)
		% \kern \kernSpacing%
		% \AND%
		% \kern \kernSpacing%
		% \mbox{\href{https://agreddy.com}{\raggedright agreddy.com}}%
	\end{minipage}
	\hfill % Flexible space between minipages
	% --- Middle Column: Website Link ---
	\begin{minipage}[t]{0.15\textwidth}
		\raggedleft % Right-align content in this minipage
		\href{https://agreddy.com}{agreddy.com}
	\end{minipage}
	\hfill % Flexible space
	% --- Right Column: QR Code ---
	\begin{minipage}[t]{0.11\textwidth}
		\raggedleft % Right-align content
		\qrcode{https://agreddy.com} % Ensure qrcode package is loaded by reddy_cv.sty
	\end{minipage}
\end{header}

% --- Vertical Space After Header ---
\vspace{\headerSpacing} % Defined in reddy_resume.sty

% ----------------------------------------------------------------------
% SUMMARY SECTION
% ----------------------------------------------------------------------

\section{SUMMARY}

\vspace{\entrySpacing} % Space before the summary content
\begin{samepage} % Keep the summary on one page if possible
	\begin{onecolentry}
		M.S. Computer Science graduate (May 2026) from Washington University in St. Louis with hands-on experience across machine learning, data engineering, and full-stack development through research, teaching, and industry internships. Skilled in building end-to-end ML pipelines, real-time data systems, and production web applications with Python, PyTorch/TensorFlow, SQL, Spark/Kafka/Airflow, and modern JavaScript frameworks. Seeking software, data, or ML engineering roles.
	\end{onecolentry}
\end{samepage}

% --- Vertical Space After Section ---
\vspace{\headerSpacing}

% ----------------------------------------------------------------------
% SKILLS SECTION
% ----------------------------------------------------------------------

\section{SKILLS}

\vspace{\entrySpacing}
\begin{skillcategory}{Programming Languages}
	Java, C++, Python, R, HTML, CSS, JavaScript/TypeScript, PHP, C\#, SQL, Rust.
\end{skillcategory}

\vspace{\entrySpacing}
\begin{skillcategory}{Frameworks \& Libraries}
	TensorFlow, PyTorch, Scikit-learn, NumPy, Pandas, Node.js, Express.js, Vue.js, React, Socket.IO, Gemini/OpenAI API.
\end{skillcategory}

\vspace{\entrySpacing}
\begin{skillcategory}{Tools \& Technologies}
	Git, Tableau, Power BI, MySQL, AWS, MongoDB, Snowflake, Apache Airflow/Kafka/Spark, Streamlit, Microsoft 365, Google Workspace.
\end{skillcategory}

\vspace{\entrySpacing}
\begin{skillcategory}{Core Competencies}
	Analytical Skills, Problem Solving, Critical Thinking, Communication, Team Collaboration, Leadership, Adaptability, Project Management.
\end{skillcategory}

\vspace{\entrySpacing}
\begin{skillcategory}{Languages}
	Proficient in English, Kannada, Telugu; Conversational in Hindi.
\end{skillcategory}

% --- Vertical Space After Section ---
\vspace{\headerSpacing}

% ----------------------------------------------------------------------
% EDUCATION SECTION
% ----------------------------------------------------------------------

\section{EDUCATION}

\vspace{\entrySpacing}
\begin{educationentry}
	{M.S. in Computer Science} % Degree
	{Washington University in St. Louis, St. Louis, Missouri, USA} % Institution & Location
	{August 2024 - May 2026} % Date
	\item \textbf{GPA:} 3.67/4.00
	\item \textbf{Thesis:} \href{https://openscholarship.wustl.edu/eng_etds/1341/}{Toward Vehicle-Agnostic Driving Signatures for Cognitive Impairment Prediction from Naturalistic Driving Data}
\end{educationentry}

\vspace{\entrySpacing}
\begin{educationentry}
	{B.A. in Computer Science and Data Analytics; Economics Minor} % Degree
	{Ohio Wesleyan University, Delaware, Ohio, USA} % Institution & Location
	{August 2020 - May 2023} % Date
	\item \textbf{GPA:} 3.42/4.00; completed the four-year program in three years
	\item \textit{Honors:} Golden Bishop Award, Mortar Board, Florence Leas Prize, Dean's List (Spring 2022, Spring 2023)
\end{educationentry}

% --- Vertical Space After Section ---
\vspace{\headerSpacing}

% ----------------------------------------------------------------------
% EXPERIENCE SECTION
% ----------------------------------------------------------------------

\section{\sectionlink{https://agreddy.com/experience/}{EXPERIENCE}}

\vspace{\entrySpacing}
\begin{experienceentry}
	{Graduate Teaching Assistant - CSE 5114: Data Manipulation and Management at Scale} % Position
	{Washington University in St. Louis, St. Louis, Missouri, USA} % Organization & Location
	{January 2026 - May 2026} % Date
	\item Supported a graduate course on large-scale data pipelines using Python, Snowflake, SQL, Airflow, Spark, Kafka, and Flink, including fault tolerance and scaling.
	\item Conducted oral midterms structured as technical interviews and mentored three capstone teams building dashboards from open data sources.
\end{experienceentry}

\vspace{\entrySpacing}
\begin{experienceentry}
	{AI Engineer Intern} % Position
	{Crittero, Inc., Remote} % Organization & Location
	{June 2025 - August 2025} % Date
	\item Built CritterCircle's (Crittero's pet platform) social media simulation, generating training data across seven personas via time-series modeling in TypeScript.
	\item Developed a TensorFlow/Keras recommendation model with a collaborative-filtering fallback and a configurable Python CLI for reproducible evaluation.
\end{experienceentry}

\vspace{\entrySpacing}
\begin{experienceentry}
	{Graduate Assistant with the Taylor Family Center for Student Success} % Position
	{Washington University in St. Louis, St. Louis, Missouri, USA} % Organization & Location
	{August 2024 - May 2026} % Date
	\item Organized engagement events for first-generation, limited-income (FLI) students (Welcome Dinner, First-Gen Week, End of Week Unwind), and managed the center's website and social media.
	\item Analyzed engagement, survey, and social media data in compliance with FERPA, developing process manuals to establish data evaluation best practices for staff.
\end{experienceentry}

\vspace{\entrySpacing}
\needspace{6\baselineskip}
\begin{experienceentry}
	{Data Analyst Intern} % Position
	{Lab714, Boca Raton, Florida, USA} % Organization & Location
	{June 2023 - May 2024} % Date
	\item Analyzed IoT environmental sensor data on AWS, building dashboards that helped clients cut pool maintenance costs by 20\%.
	\item Designed and engineered the UI for an upcoming all-in-one environmental sensor device in Figma and React.
\end{experienceentry}

% --- Vertical Space After Section ---
\vspace{\headerSpacing}

% ----------------------------------------------------------------------
% RESEARCH EXPERIENCE SECTION
% ----------------------------------------------------------------------

\section{\sectionlink{https://agreddy.com/experience/}{RESEARCH EXPERIENCE}}

\vspace{\entrySpacing}
\begin{experienceentry}
	{\href{https://openscholarship.wustl.edu/eng_etds/1341/}{Toward Vehicle-Agnostic Driving Signatures for Cognitive Impairment Prediction from Naturalistic Driving Data}} % Title
	{DRIVES Project, Washington University in St. Louis, St. Louis, Missouri, USA} % Org & Location
	{August 2025 - May 2026} % Date
	\item Built an end-to-end pipeline processing 26,968 participant-weeks of naturalistic driving data from 304 participants, deriving weekly driving features and integrating demographic covariates for CDR status prediction.
	\item Compared six CDR prediction models (GRU-DANN, DANN, logistic regression, random forest, XGBoost, MLP) under leave-one-group-out cross-validation; GRU-DANN achieved highest participant-level ROC AUC (0.599) and balanced accuracy (0.584).
\end{experienceentry}

% --- Vertical Space After Section ---
\vspace{\headerSpacing}

% ----------------------------------------------------------------------
% LEADERSHIP SECTION
% ----------------------------------------------------------------------

\section{LEADERSHIP}

\vspace{\entrySpacing}
\begin{leadershipentry}{Graduate and Professional Student Council (GPSC)}{Washington University in St. Louis}
	\begin{positionentry}{Vice President of the Graduate-Professional Council (GPC) Chamber}{May 2025 - May 2026}
		\item Managed the Council's unified Microsoft Teams workspace and website, establishing centralized communication hubs and maintaining up-to-date information on events, governance, and resources.
		\item Represented the graduate student body to university leadership and supported the GPC President in advancing strategic council priorities.
		\item Chaired the Constitution Committee (Aug 2025 - Feb 2026), coordinating stakeholder input to draft and ratify GPSC governing documents.
	\end{positionentry}
\end{leadershipentry}

% --- Vertical Space After Section ---
\vspace{\headerSpacing}

% ----------------------------------------------------------------------
% PROJECTS SECTION
% ----------------------------------------------------------------------

\section{\sectionlink{https://agreddy.com/projects/}{PROJECTS}}

\vspace{\entrySpacing}
\begin{projectentry}
	{\href{https://github.com/agopalareddy/CSE510A_Datacenter_Cooling}{Datacenter Cooling Optimization using Deep Reinforcement Learning}} % Project Title
	{August 2024 - December 2024} % Date
	\item Implemented multiple DRL algorithms (DDQN, PPO, SAC) integrated with EnergyPlus simulations for datacenter cooling optimization, achieving up to 35.8\% improvement over baseline controllers.
	\item Developed a custom Sinergym environment simulating a small datacenter with stochastic weather conditions, focusing on the underserved small-to-mid-sized facility market.
\end{projectentry}

\vspace{\entrySpacing}
\begin{projectentry}
	{\href{https://github.com/agopalareddy/CSE-5114-Project}{MLB Statcast Real-Time Data Pipeline}} % Project Title
	{September 2025 - October 2025} % Date
	\item Built a lambda-architecture pipeline for MLB Statcast pitch-level data, using Apache Airflow for batch ingestion into Snowflake and a Kafka/Spark Streaming path for live, simulated data.
	\item Developed a Streamlit dashboard with pitch-by-pitch game replay, interactive strike-zone visualizations, and pitcher/team analytics (velocity, spin rate, strikeout rate).
\end{projectentry}

\vspace{\entrySpacing}
\begin{projectentry}
	{\href{https://github.com/agopalareddy/CSE5519-Project}{Red-Blue Visual Auto Defender (VLM Security)}} % Project Title
	{September 2025 - October 2025} % Date
	\item Built an automated red-blue teaming pipeline that generates visual prompt-injection (jailbreak) attack images and evaluates them against a vision-language model (Gemma-3-4b-it) using semantic response analysis.
	\item Auto-generated deterministic, OCR-based Python defenses for each successful attack and validated them with accuracy/precision/recall metrics in an iterative refinement loop.
\end{projectentry}

\vspace{\entrySpacing}
\begin{projectentry}
	{\href{https://agreddy.com/storybook/}{Interactive Storybook}} % Project Title
	{August 2024 - December 2024} % Date
	\item Built an interactive app with Vue.js, Node.js, and MongoDB, enabling AI-driven content generation and branching story narratives.
	\item Integrated Google Gemini API and fal.ai for dynamic text/image generation, offering multiple story progression choices.
\end{projectentry}

% ----------------------------------------------------------------------
% End of Document
% ----------------------------------------------------------------------
\end{document}
